\chapter{Типи нейронів за електричною функцією}
\chaptermark{Нейрони як прилади}
\label{sec:eetypes}

У розділі~\ref{sec:neurontypes} ми розсортували нейрони за тим, який медіатор викидає їхній вихід. Електронник розсортував би їх інакше: за тим, що нейрон робить зі своїм сигналом, --- яким приладом він працює. Головний прилад книги ми знаємо з розділу~\ref{sec:glutcomputer}: інтегрувальний нейрон із витоком~\eqref{eq:glutneuron}, RC-коло з компаратором. Але в мозку є й інші. У таблиці~\ref{tab:eetypes} їх зібрано в чотири групи, і майже кожен уже траплявся в книзі.

\begin{center}
\footnotesize
\setlength{\tabcolsep}{3pt}
\renewcommand{\arraystretch}{1.15}
\begin{longtable}{>{\raggedright}p{3.1cm}>{\raggedright}p{3.3cm}>{\raggedright}p{4.7cm}>{\raggedright\arraybackslash}p{2.4cm}}
\caption{Нейрони як прилади: назва, аналог в електроніці, що прилад робить і де він трапляється в книзі.}\label{tab:eetypes}\\
\toprule
Нейрон & Прилад & Що робить & Де в книзі \\
\midrule
\endfirsthead
\toprule
Нейрон & Прилад & Що робить & Де в книзі \\
\midrule
\endhead
\bottomrule
\endfoot
\multicolumn{4}{l}{\emph{Фільтри й інтегратори}} \\
інтегрувальний нейрон із витоком & RC-інтегратор із компаратором & додає входи у вікні $\tau$ і спрацьовує біля порога & розд.~\ref{sec:glutcomputer}, \eqref{eq:glutneuron} \\
детектор збігів & І з часовим вікном & спрацьовує, лише якщо входи надійшли майже разом; дуже мале $\tau$ & розд.~\ref{sec:glutcomputer}, \ref{sec:code}, \eqref{eq:window} \\
фазичний нейрон & фільтр верхніх частот, детектор фронту & відповідає на зміну входу й замовкає & розд.~\ref{sec:sensors} \\
адаптивний нейрон & фільтр верхніх частот з автоматичним регулюванням підсилення & частота падає за сталого входу & розд.~\ref{sec:sensors}, \eqref{eq:nakarushton} \\
резонатор & коливальний контур, смуговий фільтр & найкраще відповідає на вхід своєї частоти & підрозділ~\ref{sec:resonator} \\
інтегратор без витоку & інтегратор на операційному підсилювачі & перетворює швидкість на положення й тримає його & підрозділ~\ref{sec:perfectint} \\
\midrule
\multicolumn{4}{l}{\emph{Перетворювачі}} \\
тонічний нейрон & перетворювач напруга---частота & частота лінійно йде за входом і майже не втомлюється & розд.~\ref{sec:code} \\
градуальний нейрон (без імпульсів) & аналоговий підсилювач, буфер & передає плавну напругу на коротку відстань & розд.~\ref{sec:sensors} \\
випрямний нейрон & однопівперіодний випрямляч & частота не буває від'ємною, $[u]_+$ & розд.~\ref{sec:glutcomputer}, \ref{sec:gaba} \\
пара «увімкнення---вимкнення» & пара випрямлячів, двопровідний код & один відповідає на «більше», другий --- на «менше» & розд.~\ref{sec:glutcomputer}, \eqref{eq:dualrail} \\
нейрон із шунтувальним гальмуванням & аналоговий дільник, регулювання підсилення & ділить вхід, а не віднімає від нього & розд.~\ref{sec:gaba}, \eqref{eq:shunt} \\
\midrule
\multicolumn{4}{l}{\emph{Генератори й час}} \\
ритмоводій & релаксаційний генератор, мультивібратор & сам видає імпульси зі сталою частотою & розд.~\ref{sec:serocomputer}, \ref{sec:dofa} \\
ритмоводій із входом & генератор, керований напругою & власний ритм, який вхід прискорює чи сповільнює & розд.~\ref{sec:chemoutput} \\
пачковий нейрон & стробований генератор пачок & видає пачки частих імпульсів --- «тире» & розд.~\ref{sec:code}, \eqref{eq:burst} \\
нейрон, прив'язаний до фази & фазовий детектор, фазове автопідстроювання & спрацьовує на певній фазі вхідної хвилі & розд.~\ref{sec:sensors}, \ref{sec:chemoutput} \\
аксон із затримкою & лінія затримки & затримує сигнал на заданий час & розд.~\ref{sec:glutcomputer}, рис.~\ref{fig:itd} \\
\midrule
\multicolumn{4}{l}{\emph{Пам'ять і перемикання}} \\
бістабільний нейрон & тригер Шмітта, засувка & короткий вхід переводить його в тривалий стан & підрозділ~\ref{sec:bistable} \\
нейрон із дендритними гілками & мультиплексор, маленька багатошарова мережа & гілка пропускає вхід, лише якщо є дозвільний вхід & розд.~\ref{sec:synthesis} \\
\end{longtable}
\end{center}

Одне застереження важливіше за всю таблицю: це не різні клітини, а різні \emph{режими}. Той самий нейрон може бути інтегратором за повільного входу й детектором збігів, коли гальмування вкоротило його $\tau$ (розділ~\ref{sec:code}); ритмоводієм під час неспання й мовчазним приймачем уві сні. Який прилад вийде, вирішують іонні канали мембрани й широкомовні канали, що їх підлаштовують.

\section{Чотири прилади, яких у книзі ще не було}

\subsection{Резонатор}
\label{sec:resonator}

Деякі нейрони мають повільний струм, який протидіє будь-якій зміні напруги: напруга зростає --- струм тягне її вниз, падає --- тягне вгору, але із запізненням. Для електронника такий струм поводиться як індуктивність, а разом з ємністю мембрани вони утворюють коливальний контур. Нейрон найсильніше відповідає на вхід певної частоти, зазвичай від одиниць до одного-двох десятків герц, і слабше --- на повільніший і швидший~\cite{HutcheonYarom2000}. Це смуговий фільтр в одній клітині: із шумного входу він вибирає ритм своєї частоти, наприклад місцевий ритм розділу~\ref{sec:gaba}.

\subsection{Інтегратор без витоку}
\label{sec:perfectint}

Інтегратор із витоком пам'ятає вхід лише час $\tau$ --- десятки мілісекунд. Але окові, щоб стояти нерухомо, треба пам'ятати своє положення секундами: команда «повернути» надходить як короткочасна швидкість, а тримати око треба в новому положенні. Мережа розв'язує це тим самим зворотним зв'язком, що й пам'ять розділу~\ref{sec:glutcomputer}. Якщо група нейронів зі сталою часу $\tau$ збуджує сама себе з вагою $w$, то $\tau\,\dd r/\dd t=-r+w\,r+I$, і стала часу групи дорівнює
\begin{empheq}[box=\keybox]{equation}
\tau_{\text{еф}} \;=\; \frac{\tau}{1-w} .
\label{eq:perfectint}
\end{empheq}
За $\tau=0.1\,\text{с}$ і $w=0.995$ виходить $20\,\text{с}$: майже ідеальний інтегратор із нейронів, кожен з яких сам по собі забуває за десяту частку секунди~\cite{Seung1996}. Ціна --- точне налаштування: коли $w$ трохи більше за одиницю, інтегратор іде в насичення, коли трохи менше --- швидко забуває. Тому такому інтеграторові потрібне постійне підлаштування навчанням, на кшталт того, що описано в розділі~\ref{sec:motorlearning}.

\subsection{Бістабільний нейрон}
\label{sec:bistable}

У розділах~\ref{sec:glutcomputer} і~\ref{sec:gaba} тригер збирався з групи нейронів. Але буває тригер і в одній клітині. Деякі нейрони мають струм, який, щойно ввімкнувшись, сам підтримує збудження: короткий вхід переводить нейрон у тривалу роботу --- \emph{плато}, --- а гальмування повертає в спокій. Це тригер Шмітта: у нейрона два стійкі стани й гістерезис між ними. Так поводяться, наприклад, \nt{ACET}-нейрони, що керують м'язами, і вмикається ця властивість широкомовним каналом: плато з'являється, коли до нейрона доходить серотонін~\cite{Hounsgaard1988}. Та сама клітина із \nt{SERO} --- засувка, без нього --- звичайний інтегратор.

\subsection{Інтегратори й резонатори}

Теорія поділяє збудливі клітини на два класи~\cite{Izhikevich2007}. \emph{Інтегратори} схожі на генератор, керований напругою, який починає з нуля: трохи вище порога нейрон працює як завгодно повільно, і частота плавно зростає зі входом. Такий нейрон додає все, що надійшло, і добре передає частотою величину входу. \emph{Резонатори} схожі на генератор із мінімальною частотою: повільно працювати вони не вміють, за порогового входу одразу видають імпульси зі скінченною частотою і віддають перевагу входу своєї частоти. Перші --- природні прилади частотного коду розділу~\ref{sec:code}, другі --- часового.

\begin{keyblock}
\begin{proposition}[Один нейрон --- багато приладів]
\label{prop:eetypes}
За електричною функцією нейрони працюють як інтегратори з витоком і без нього, детектори збігів, фільтри верхніх частот і смугові фільтри, перетворювачі напруга---частота, випрямлячі, дільники, генератори, лінії затримки й тригери. Який прилад вийде з даної клітини, задають її іонні канали й широкомовні канали, що їх підлаштовують. Самі режими виміряно; те, наскільки мозок використовує таке переналаштування для обчислень, --- здебільшого гіпотеза.
\end{proposition}
\end{keyblock}

Для інженера це схоже на програмовану логічну матрицю: залізо одне, а прилад задає конфігурація. Тільки конфігурацію тут змінюють не під час прошивання, а на ходу, --- повільними широкомовними каналами, про які йшлося в розділах~\ref{sec:serocomputer}--\ref{sec:acet}.

\section{Задачі}

\begin{exercise}[\lvlA]
\label{ex:18:1}
\emph{Допуск інтегратора без витоку.} Нейрони з $\tau=0.1$~с тримають положення ока за~\eqref{eq:perfectint} протягом $T=1$~с із похибкою не більше $5\%$ в обидва боки. (а)~Знайдіть допустимий діапазон $w$. (б)~$w$ --- сума ваг $K$ синапсів, кожен із незалежною похибкою $10\%$. За якого $K$ розкид суми вкладається в допуск?
\end{exercise}

\begin{exercise}[\lvlB]
\label{ex:18:2}
\emph{Резонатор як контур.} Повільний струм підрозділу~\ref{sec:resonator} опишемо змінною $W$: $C\,\dd V/\dd t=-g_LV-g_wW+I$, $\tau_w\,\dd W/\dd t=V-W$. Покажіть, що це $RC$-коло з паралельною гілкою $R_w=1/g_w$, $L_w=\tau_w/g_w$, і за $C/g_L=10\ms$, $\tau_w=100\ms$, $\alpha=g_w/g_L=4$ знайдіть частоту резонансу й виграш $|Z|$ у ньому над $|Z(0)|$.
\end{exercise}

\begin{exercise}[\lvlB]
\label{ex:18:3}
\emph{Як інтегратор починає працювати.} (а)~Нейрон $\tau\,\dd V/\dd t=-V+\mu$, $\tau=10\ms$, з порогом $\theta$ і скиданням у нуль. За якого $\mu/\theta-1$ частота дорівнює $1$~Гц? (б)~Яку частоту дає модель $\tau\,\dd v/\dd t=v^2+\mu$ (імпульс --- відхід $v$ у $+\infty$, скидання в $-\infty$) за $\mu=0.01$?
\end{exercise}

\begin{exercise}[\lvlB]
\label{ex:18:4}
\emph{Моделювання: інтегратор і резонатор.} Модель задачі~\ref{ex:18:2} з $g_L=1$, $\tau_m=10\ms$, $\tau_w=100\ms$, порогом $1$ і скиданням $V\to0$ ($W$ не змінюється) отримує $\mu=1.5\sin2\pi ft$, $f=1$--$40$~Гц. У якій смузі частот нейрон видає імпульси за $\alpha=0$ і за $\alpha=4$? Порівняйте з умовою $1.5\,|Z(f)|>1$.
\end{exercise}

\begin{exercise}[\lvlB]
\label{ex:18:5}
\emph{Засувка з однієї клітини.} Нейрон зі струмом плато: $\tau\,\dd r/\dd t=-r+F(wr+I)$, $F(x)=\min(1,\max(0,x))$, $\tau=10\ms$, $w=1.5$, фон $I=-0.2$. (а)~Якої найменшої тривалості імпульс $I=0.3$ переводить нейрон з $r=0$ у плато? (б)~Якої найменшої амплітуди $B$ довгий імпульс $I=-0.2-B$ повертає його в спокій?
\end{exercise}

\begin{exercise}[\lvlB]
\label{ex:18:6}
\emph{Самоналаштування інтегратора.} Під час фіксації погляду вхід $I=0$, датчик ковзання дає швидкість відходу $e=\dd r/\dd t$, і $\dd w/\dd t=-\eta\,e\,r$. За $w=0.95$, $\tau=0.1$~с, $\langle r^2\rangle=1$ знайдіть $\eta$, за якого за $60$~с фіксацій $|1-w|$ увійде в допуск задачі~\ref{ex:18:1}. Що зламається, якщо вчитися й під час поворотів ока?
\end{exercise}

\begin{exercise}[\lvlC]
\label{ex:18:7}
\emph{Навіщо переналаштовувати прилад?} Широкомовний канал перемикає $\alpha$ моделі задачі~\ref{ex:18:2} між $0$ і $4$. У скільки разів резонатор виграє в інтегратора у відношенні сигнал/шум для слабкого синуса в білому шумі струму за $11$~Гц і за $1$~Гц? Придумайте задачу, де вигідне саме перемикання режиму.
\end{exercise}
